\documentclass[a4paper,12pt]{jsarticle}
\usepackage[margin=25mm]{geometry}
\usepackage{amsmath,amssymb}
\usepackage{enumitem}

\begin{document}

\begin{center}
{\LARGE 数I・A 共通テスト本番レベル（やや難）}
\end{center}

\vspace{1em}

\section*{問題}

\section*{第1問（数と式・二次関数）}
実数 $a$ を定数とし、
\[
f(x)=x^2-2ax+a+2
\]
とする。
\begin{enumerate}[label=\arabic*.]
  \item $f(x)$ の最小値を $a$ を用いて表せ。
  \item 不等式 $f(x)\le 0$ の解集合の長さが $2$ 以上となるような $a$ の範囲を求めよ。
  \item 方程式 $f(x)=0$ が $0<\alpha<1<\beta$ を満たす2つの実数解 $\alpha,\beta$ をもつような $a$ の範囲を求めよ。
  \item $a=4$ のとき、不等式 $|f(x)|\le 3$ を満たす $x$ の範囲を求めよ。
\end{enumerate}

\section*{第2問（場合の数・確率）}
1から12までの整数が1枚ずつ書かれた12枚のカードから、同時に4枚を取り出す。
\begin{enumerate}[label=\arabic*.]
  \item 全事象の数を求めよ。
  \item 4枚のうち偶数がちょうど2枚である確率を求めよ。
  \item 「最大値が10である」という条件のもとで、最小値が3である条件付き確率を求めよ。
  \item 4枚の和が3の倍数となる確率を求めよ。
\end{enumerate}

\section*{第3問（図形と計量）}
三角形 $ABC$ において、
\[
AB=6,\quad AC=10,\quad \angle A=120^\circ
\]
とする。辺 $BC$ の中点を $M$ とする。
\begin{enumerate}[label=\arabic*.]
  \item 辺 $BC$ の長さを求めよ。
  \item 三角形 $ABC$ の面積を求めよ。
  \item 三角形 $ABC$ の内接円の半径 $r$ を求めよ。
  \item 中線 $AM$ の長さを求めよ。
\end{enumerate}

\section*{第4問（データの分析）}
ある20人の得点分布が次の表で与えられる。
\[
\begin{array}{c|cccccccc}
\text{得点} & 40 & 45 & 50 & 55 & 60 & 65 & 70 & 75 \\ \hline
\text{人数} & 1 & 2 & 3 & 4 & 4 & 3 & 2 & 1
\end{array}
\]
\begin{enumerate}[label=\arabic*.]
  \item 中央値（メジアン）を求めよ。
  \item 第1四分位数 $Q_1$、第3四分位数 $Q_3$、四分位範囲 $IQR$ を求めよ。
  \item 平均値 $\bar{x}$ を求めよ。
  \item 分散 $s^2=\dfrac{1}{20}\sum (x-\bar{x})^2$ と標準偏差 $s$ を求めよ。
\end{enumerate}

\end{document}
